\section{Results}
\label{sec:results}
% One subsection per figure, in the order of the argument:
%  3.1 the two regimes (Fig. 1, Table 1)   3.2 the intensity clause (Fig. 2)
%  3.3 the nitrile (Fig. 3)                3.4 benzene patterns (Fig. 4)
%  3.5 rule-out value (LR-)                3.6 replication (Fig. 5)

After quality control and the one-spectrum-per-structure rule, the confirmatory Chemotion set contained 1,539 spectra in the main stratum and the NIST set 5,124. Of the 101 evaluable rules, 73 had at least ten positive spectra in Chemotion and 85 in NIST; the remaining rules are reported in Table~\ref{tab:rules} but not interpreted. Unless stated otherwise, all figures below refer to the primary band-detection threshold and to the full rule definition, that is, position together with any intensity or shape qualifier the source gives.

\subsection{Two regimes of the spectrum}
\label{sec:regimes}

Figure~\ref{fig:regions} plots the positive likelihood ratio of every powered rule against the centre of its published window, and Table~\ref{tab:summary} counts the rules in each \LRp{} band by region. The two collections give the same picture. Rules whose window lies at or above 1500~\wn{} have a median \LRp{} of 2.8 in Chemotion and 3.0 in NIST; rules whose window lies at or below 1350~\wn{} have a median of 1.06 and 1.21. In the upper region 7 of 31 Chemotion rules and 9 of 37 NIST rules reach \LRp{} 5, the level that the diagnostic literature uses, by convention, to mark a test whose positive result carries substantial weight on its own~\cite{deeks2004likelihood}; in the fingerprint region none of the 38 or 43 rules does, and 26 of 38 in Chemotion (21 of 43 in NIST) fall below 1.2, where the band is found almost as often without the group as with it. The handful of rules whose window straddles the boundary (4 and 5) behave like the fingerprint.

The rules that carry most evidence are the ones a spectroscopist would name first. In Chemotion the seven rules above \LRp{} 5 are the saturated ester carbonyl (43), the aromatic nitrile (189, but see Section~\ref{sec:nitrile}), the two asymmetric NH$_2$ stretches of primary amines (8.7 alkyl, 6.8 aryl), the aromatic ester and saturated ketone carbonyls (6.1 and 5.9) and the asymmetric nitro stretch (5.3). NIST, whose composition gives power to groups that Chemotion lacks, adds the three organic carbonate carbonyls (5.7 to 9.6), the saturated aldehyde carbonyl (6.3) and the terminal alkyne C$\equiv$C (12), and keeps the saturated ester carbonyl (12) and the nitriles at the top. Below 5 but above 2 sit most of the remaining carbonyl, N-H, O-H and C$\equiv$C rules, together with the C-H stretches of methyl and methylene groups.

The fingerprint tells a different story. Its rules are not failing to find the band: the median sensitivity of the 38 Chemotion fingerprint rules is 0.90, higher than the 0.47 of the rules above 1500~\wn{}. What they lack is specificity, a median of 0.16 against 0.85, and for 15 of the 38 it is below 0.10. The C-O stretch of secondary alcohols (1075 to 1150~\wn{}) illustrates the pattern: 94\% of secondary alcohols show a band in the window, and so do 95\% of the compounds that are not secondary alcohols, for an \LRp{} of 0.98. The same holds for the C-O of phenols, ethers and acids, the C-N of amines and amides, and the out-of-plane C-H bends of alkenes and benzene rings. The region is simply crowded: at the primary threshold a typical Chemotion spectrum carries 16 detected bands between 700 and 1350~\wn{} (interquartile range 13 to 18), enough to occupy 85\% of the 50~\wn{} windows in that range, against 10 bands in the whole region above 1500~\wn{}. A rule that asks only whether a band is present in a fingerprint window therefore has little room to discriminate.

NIST is somewhat kinder to the fingerprint (median specificity 0.34; 7 of 43 rules below 0.10), and nine of its fingerprint rules reach \LRp{} 2, against four in Chemotion. The difference is composition rather than technique: a third of the NIST compounds carry no aromatic ring, and their fingerprint is emptier, 12 bands per spectrum and 69\% of the windows occupied. The rules that gain most are those that pair a position with a second condition, the monosubstituted benzene ring bend together with its wag (3.7) and the carbonate O-C-O stretch (3.7), and the vinyl C-H wag at 990~\wn{} (4.0), whose window is narrow. None of them reaches 5. The conclusion is the same in both collections: at or below 1350~\wn{}, band position alone is close to uninformative, and whatever diagnostic value the region holds has to come from intensity, shape, pattern or combination, which Sections~\ref{sec:intensity} to \ref{sec:benzene} examine.

\begin{figure}[t]
\centering
\includegraphics[width=\textwidth]{fig1_lr_by_region}
\caption{Positive likelihood ratio of every powered rule (at least 10 spectra containing the group) against the centre of its published window, in the Chemotion (top) and NIST (bottom) collections. Whiskers are 95\% scaffold-cluster bootstrap intervals; the shaded area marks windows at or below 1350~\wn{}. Median \LRp{} is 2.8 and 3.0 for rules at or above 1500~\wn{}, 1.1 and 1.2 in the fingerprint.}
\label{fig:regions}
\end{figure}

% generated by scripts/tables.py, do not edit by hand
\begin{table}[t]\centering\small
\caption{Number of powered rules in each LR$^+$ band, by spectral region. LR$^+ \geq 5$ is the usual threshold for a test that changes a decision on its own; LR$^+ < 1.2$ means the band is present about as often without the group as with it.}\label{tab:summary}
\begin{tabular}{@{}lrrrrrr@{}}\toprule
Region & $\geq 5$ & 2 to 5 & 1.2 to 2 & $<1.2$ & rules & median LR$^+$ \\
\midrule
\multicolumn{7}{@{}l}{\emph{All powered rules}} \\
\quad Chemotion & 7 & 23 & 14 & 29 & 73 & 1.47 \\
\quad NIST & 10 & 31 & 22 & 22 & 85 & 1.85 \\
\midrule
\multicolumn{7}{@{}l}{\emph{Window at or above 1500 cm$^{-1}$}} \\
\quad Chemotion & 7 & 17 & 5 & 2 & 31 & 2.80 \\
\quad NIST & 9 & 22 & 5 & 1 & 37 & 3.03 \\
\midrule
\multicolumn{7}{@{}l}{\emph{Window straddling 1350--1500 cm$^{-1}$}} \\
\quad Chemotion & 0 & 2 & 1 & 1 & 4 & 1.75 \\
\quad NIST & 1 & 0 & 4 & 0 & 5 & 1.40 \\
\midrule
\multicolumn{7}{@{}l}{\emph{Window at or below 1350 cm$^{-1}$ (fingerprint)}} \\
\quad Chemotion & 0 & 4 & 8 & 26 & 38 & 1.06 \\
\quad NIST & 0 & 9 & 13 & 21 & 43 & 1.21 \\
\bottomrule
\end{tabular}
\end{table}


\subsection{What the intensity clause buys and costs}
\label{sec:intensity}

Forty-nine of the powered rules state an intensity, ``strong'' in all but the nitro scissors (``medium''). Figure~\ref{fig:intensity} compares, for each of them, the \LRp{} obtained from position alone with the \LRp{} obtained when the band must also reach the stated intensity class, defined here as at least 0.6 (strong) or 0.2 (medium) of the strongest band in the spectrum. The clause does what it is meant to do: median specificity rises from 0.39 to 0.85 in Chemotion and from 0.59 to 0.88 in NIST, and \LRp{} improves by more than 20\% in 16 of 20 and 23 of 29 rules. It also does what the tables never mention: median sensitivity falls from 0.78 to 0.44 and from 0.70 to 0.39, and in 8 rules of each collection it is halved or worse.

Whether the trade pays depends on how strong the band really is. Where the group gives one of the strongest bands in the spectrum, the clause costs little and buys a lot: the saturated ester carbonyl goes from 2.9 to 43 while keeping 40 of its 50 percentage points of sensitivity, the aromatic acid carbonyl from 2.4 to 4.8 (0.69 to 0.62), the saturated ketone carbonyl from 3.0 to 5.9 (0.66 to 0.58). Two fingerprint rules are rescued this way: the symmetric nitro stretch at 1330 to 1390~\wn{} and the aromatic ester C-C-O stretch at 1250 to 1310~\wn{} are worthless by position (1.1 in both collections) and become useful with the clause (4.7 and 4.4 in Chemotion, 5.1 and 3.4 in NIST), because a nitro group or an aryl ester does put one of the strongest bands of the spectrum there, and most other compounds do not. Where the band is strong only in the author's memory, the clause removes most true positives: the acid O-H stretch at 2500 to 3500~\wn{} keeps its specificity gain but falls from sensitivity 1.00 to 0.10 (0.44 in NIST), the secondary amide N-H bend from 0.94 to 0.44, the aromatic nitrile from 0.78 to 0.06 (Section~\ref{sec:nitrile}). In four Chemotion rules and one NIST rule the clause lowers \LRp{} outright, which happens when the qualifier removes true positives faster than false ones.

The clause has a second cost that Figure~\ref{fig:intensity} does not show. A rule that misses most true positives loses its value for ruling a group out: in NIST the median \LRn{} of these rules worsens from 0.55 with position alone to 0.72 with the clause, and the number of rules whose \LRn{} exceeds 0.8, that is, whose absence says almost nothing, rises from 4 to 11 of 29. The qualifier turns a rule that worked in both directions into one that works only when the band happens to be strong.

\begin{figure}[t]
\centering
\includegraphics[width=\textwidth]{fig2_position_vs_full}
\caption{Rules that state an intensity: \LRp{} with position alone (open circle) and with the intensity clause applied (filled), in Chemotion (left, 20 rules) and NIST (right, 29 rules). The right-hand column gives the sensitivity before and after the clause. The clause raises \LRp{} in most rules but halves sensitivity or worse in 8 of 20 and 8 of 29.}
\label{fig:intensity}
\end{figure}

\subsection{The nitrile band}
\label{sec:nitrile}

The C$\equiv$N stretch is the clearest case, because the band itself is not in doubt: it is sharp (median width at half height 14~\wn{} in Chemotion, 19 in NIST), it sits in a region of the spectrum where almost nothing else absorbs, and it is detected at the primary threshold in 68 of the 72 Chemotion nitriles and 115 of the 127 NIST nitriles. Yet the published rules, which describe it as ``intense'', recognise 6\% of the aromatic nitriles in Chemotion and 3\% and 6\% of the saturated and aromatic nitriles in NIST.

Figure~\ref{fig:nitrile}a shows why. For each nitrile the height of the C$\equiv$N band relative to the strongest band of its spectrum is a median of 0.30 in Chemotion (interquartile range 0.08 to 0.50) and 0.18 in NIST (0.08 to 0.36); it reaches the ``strong'' class, 0.6 of the strongest band, in one nitrile in ten in either collection (10\% and 11\%). The stretch of a triple bond between two atoms of similar electronegativity carries a modest dipole change, and in a molecule that also has a carbonyl, an ether or an aromatic ring, one of those will provide the strongest band. In Chemotion 19 of the 72 nitriles have the band below 5\% of the maximum; most belong to one family of carbazolyl dicyanobenzenes whose spectra are dominated by the aromatic bands. In NIST 28 of 127 fall below the same level, and they are small, unrelated molecules. Transmission spectra of KBr discs and mulls give the same distribution as ATR spectra of neat solids, so the weakness of the band is a property of the vibration, not of the sampling technique.

Position is a second, smaller problem (Figure~\ref{fig:nitrile}b). The source gives two 20~\wn{} windows, 2220 to 2240 for aromatic and 2240 to 2260 for saturated nitriles. In Chemotion, where the aromatic nitriles come from a few closely related series, 82\% of the apexes fall inside the aromatic window. In NIST they spread from 2190 to 2280~\wn{}, and the windows catch 47\% of the aromatic and 45\% of the saturated nitriles: conjugation, $\alpha$-substitution and the solid state move the band by more than the windows allow. By position alone the rules therefore reach a sensitivity of 0.78 in Chemotion but 0.46 and 0.39 in NIST, with a specificity of 0.92 to 0.99 and an \LRp{} of 10 to 32; the intensity clause then removes what is left.

The numbers that come out of the full rule are a warning about reading \LRp{} alone. With the clause the specificity is 1.00, no non-nitrile ever shows a strong band there, and the \LRp{} is 189 in Chemotion and 62 and 161 in NIST, the highest values in the whole table. A reader who looked only at that column would call the nitrile rule the best in the book. It is, in the sense that a strong band at 2230~\wn{} is proof of a nitrile; it is also a rule that will miss nine nitriles in ten, with an \LRn{} of 0.94 to 0.97 that makes its absence meaningless. What the source means by ``intense'' is, we think, ``conspicuous'': narrow, isolated and impossible to overlook. That is a statement about how easy the band is to notice, not about its height, and the two should not be confused when a rule is written down.

\begin{figure}[t]
\centering
\includegraphics[width=\textwidth]{fig3_nitrile}
\caption{The C$\equiv$N stretch in 72 Chemotion and 127 NIST nitriles. (a) Fraction of nitriles whose band, detected between 2190 and 2280~\wn{}, reaches a given height relative to the strongest band of the spectrum; the dashed lines mark the "medium" and "strong" cut points. One nitrile in ten reaches "strong" in either collection. (b) Apex position by class against the two published windows: the aromatic window catches 82\% of Chemotion but 47\% of NIST aromatic nitriles, the saturated window 45\% of NIST saturated nitriles.}
\label{fig:nitrile}
\end{figure}

\subsection{Benzene substitution patterns}
\label{sec:benzene}

The out-of-plane C-H bending rules for benzene substitution patterns deserve their own figure for two reasons: they are the rules students are taught to trust in the fingerprint, and they are the only rules in the set that combine two bands or use the absence of a band as evidence. Figure~\ref{fig:benzene} gives their \LRp{} and \LRn{} in both collections. Chemotion, with 93\% aromatic compounds and many molecules carrying two or three rings, is the hardest environment these rules could meet; NIST, with 66\% aromatic compounds and mostly one ring per molecule, is closer to the spectra the rules were written from.

As rule-in evidence none of them reaches \LRp{} 5, and most stay below 2. The monosubstituted wag at 710 to 770~\wn{} is found in 97\% of monosubstituted rings and in 93\% of everything else (\LRp{} 1.05 in Chemotion, 1.24 in NIST). The ring bend at 680 to 700~\wn{} does better, 1.9 and 3.1, and asking for both bands adds little (2.0 and 3.7). Every mono and meta rule scores higher in NIST than in Chemotion, and the reason is the denominator: in NIST a third of the compounds have no ring to put a band at 690~\wn{}, whereas in Chemotion 48\% of the compounds that are not monosubstituted still show one.

The two absence rules are the weakest in the table. The ortho rule (wag at 735 to 770, no band at 680 to 700~\wn{}) gives \LRp{} 0.82 in Chemotion, below 1: an ortho-disubstituted compound in this collection is less likely to satisfy the rule than a compound that is not. The reason is that 74\% of the ortho compounds do show a band at 690~\wn{}, against 56\% of the others, because the ortho ring usually carries a monosubstituted phenyl elsewhere in the molecule. In NIST, with simpler molecules, the rule recovers to 1.7 (47\% against 39\%). The para rule behaves the same way (1.2 and 1.1). A rule written for a molecule with one ring turns against itself when the molecule has three.

As rule-out evidence the picture changes. The absence of the 690~\wn{} bend lowers the odds of a monosubstituted ring six-fold (\LRn{} 0.17 and 0.15), and the absence of the meta wag does the same for a meta ring (0.25 and 0.14), because these bands are almost always present when the pattern is (sensitivity 0.89 to 0.98). The ortho and para rules, with \LRn{} 0.7 to 1.1, are uninformative in both directions. The substitution rules, then, do not identify a pattern; when the expected band is missing, they exclude one, and that is the use a spectroscopist can still make of them.

\begin{figure}[t]
\centering
\includegraphics[width=\textwidth]{fig4_benzene}
\caption{Rule-in (\LRp{}) and rule-out (\LRn{}) value of the out-of-plane bending rules for benzene substitution patterns, with 95\% intervals, in both collections. No pattern reaches \LRp{} 5; the ring bend at 690~\wn{} and its combination with the wag are the only rules above 2, and only in NIST. The monosubstituted ring bend is a useful rule-out (\LRn{} 0.17 and 0.15).}
\label{fig:benzene}
\end{figure}

\subsection{Rule-out value}
\label{sec:ruleout}

The tables read in one direction, from band to group; the likelihood ratios read in both. Of the 73 and 85 powered rules, 11 and 17 have an \LRn{} of 0.3 or lower, that is, their absence lowers the odds of the group at least three-fold, and 25 and 24 have an \LRn{} of 0.8 or higher, so that their absence means almost nothing. Only one rule, the terminal alkyne C-H stretch at 3250 to 3350~\wn{}, is useful in both directions in both collections (\LRp{} 3.8 and 4.5, \LRn{} 0.17 and 0.06): the band is present in 87\% and 96\% of terminal alkynes and in about a fifth of everything else. The other reliable rule-outs are the secondary amide N-H stretch (\LRn{} 0.08 in Chemotion; 0.33 in NIST, where hydrogen-bonded solids broaden it), the primary amide N-H stretches (0.10 in NIST), the primary amine NH$_2$ scissors (0.11), the monosubstituted benzene ring bend (0.17 and 0.15) and the meta wag (0.25 and 0.14).

An \LRn{} is only usable when the negative result actually occurs. Several fingerprint rules have an \LRn{} below 0.2 with an \LRp{} of about 1, for instance the tertiary alcohol C-O stretch at 1100 to 1210~\wn{} in NIST (\LRn{} 0.08, \LRp{} 1.07), and this is not a contradiction: the band is present in every tertiary alcohol and in 93\% of the other compounds, so its absence does rule the group out, but the absence is seen in one spectrum in fourteen. These rules are listed in Table~\ref{tab:rules} with their specificity, which is the number that tells the reader how often the rule-out can be invoked.

\subsection{Agreement between collections}
\label{sec:replication}

Everything above could be a property of one laboratory, one instrument and one sampling technique. Figure~\ref{fig:replication} plots, for the 71 rules powered in both collections, the \LRp{} in Chemotion against the \LRp{} in NIST. The rank correlation is 0.79, and 62 of the 71 rules (87\%) agree within a factor of two, with the fingerprint rules clustered around (1, 1) and the carbonyl, N-H and C$\equiv$X rules rising along the diagonal. Two collections that share no compound, no instrument, no sampling technique and no decade order the rules the same way.

The nine rules that disagree by more than a factor of two are labelled, and each has a reading. Six involve groups that are rare or narrowly represented in Chemotion: the saturated primary and secondary amines (13 and 24 positives), the saturated acid (10) and the aromatic aldehyde (64, with a wide interval), where the Chemotion estimate is simply uncertain, and the two primary amine NH$_2$ stretches (13 and 61 positives), which score higher in Chemotion than in NIST in the direction expected from hydrogen bonding: in KBr discs and mulls the N-H bands of amines broaden and shift down, out of the narrow windows. Two are chemistry of the collection rather than of the rule: the saturated ester carbonyl falls from 43 to 12 because NIST holds many more anhydrides, lactones and acids that place a strong band in the same window, and the terminal alkyne C$\equiv$C rises from 2.7 to 12 because the Chemotion alkynes are mostly conjugated aryl alkynes whose band sits above the 2100 to 2140~\wn{} window. The ninth is the ortho absence rule of Section~\ref{sec:benzene}. No rule changes sign: what is evidence in one collection is evidence in the other, and what is noise is noise.

Agreement is weaker for \LRn{} (72\% within a factor of two), because \LRn{} depends on sensitivity, and sensitivity depends on which members of a group a collection happens to contain: the secondary amide N-H stretch is found in 94\% of Chemotion amides but 75\% of NIST amides, most of which are hydrogen-bonded solids. Rule-in value travels between collections better than rule-out value, and the rule-out figures in Table~\ref{tab:rules} should be read with that in mind.

\begin{figure}[t]
\centering
\includegraphics[width=0.7\textwidth]{fig5_replication}
\caption{\LRp{} of the 71 rules powered in both collections, Chemotion against NIST. The shaded band is agreement within a factor of two (62 of 71 rules, 87\%); Spearman correlation 0.79. Labelled rules disagree by more than a factor of two; six of the nine involve a group with fewer than 30 positives in Chemotion.}
\label{fig:replication}
\end{figure}
